\documentclass[11pt]{article}

\usepackage[utf8]{inputenc}
\usepackage[T1]{fontenc}
\usepackage{lmodern}
\usepackage{geometry}
\usepackage{amsmath,amssymb,amsfonts,amsthm,mathtools}
\usepackage{bm}
\usepackage{enumitem}
\usepackage{microtype}
\usepackage{hyperref}
\usepackage{titlesec}
\usepackage{booktabs}
\usepackage{array}
\usepackage{setspace}
\usepackage{mathrsfs}
\usepackage{physics}

\geometry{margin=1in}
\hypersetup{colorlinks=true, linkcolor=black, urlcolor=blue, citecolor=black}
\setlist[enumerate]{topsep=0.4em,itemsep=0.35em}
\setlist[itemize]{topsep=0.3em,itemsep=0.25em}
\titleformat{\section}{\large\bfseries}{\thesection.}{0.5em}{}
\titleformat{\subsection}{\normalsize\bfseries}{\thesubsection.}{0.5em}{}
\setstretch{1.06}

\newcommand{\Aether}{\AE{}ther}
\newcommand{\Aflow}{\AE{}ther-flow}
\newcommand{\TheoryName}{\Aflow{} Interpretation of Relativity}
\newcommand{\TheoryProgram}{\Aether{} / \Aflow{} framework}
\newcommand{\CompletedMetric}{g^{\sharp}}

\newcommand{\PrePreProtoSectionPullbackCore}{\mathfrak S^{\mathrm{sub,proto,pppresec,pb}}}
\newcommand{\PullbackOperatorCore}{\mathfrak S^{\mathrm{sub,proto,pppresec,pb,op}}}
\newcommand{\PullbackBodyOperatorCore}{\mathfrak S^{\mathrm{sub,proto,pppresec,pb,bodyop}}}
\newcommand{\PullbackBodyResponseCore}{\mathfrak S^{\mathrm{sub,proto,pppresec,pb,bodyresp}}}

\newtheorem{definition}{Definition}
\newtheorem{proposition}{Proposition}
\newtheorem{remark}{Remark}
\newtheorem{corollary}{Corollary}
\newtheorem{theorem}{Theorem}

\title{\textbf{The \TheoryName{}}\\[0.4em]
\large Primitive-Reservoir Observer-Reduction\\
\large Section-Centered Hidden-Response\\
\large Pullback Body-Response Core Necessity / Uniqueness\\
\large Next Benchmark Criterion}
\author{}
\date{}

\begin{document}

\maketitle

\begin{abstract}
The current active-primary theorem already proved that the observer-relevant pullback body-response core \(\PullbackBodyResponseCore\) is the unique minimal benchmark-facing datum on the recorded section-centered hidden-response route up to body-response-faithful equivalence. Another same-output deeper-origin relay that merely preserves that same body-response core and therefore merely reproduces the same pullback body-operator core, the same pullback-operator core, the same pullback-quadratic core, the same hidden-response package, the same residual-normal-form datum, and the same benchmark-faithful pre-proto-section package no longer counts as the honest default next benchmark-facing manuscript.

This note records the routing consequence of that gain. The next honest benchmark-facing move must now derive or justify the current body-response core itself, or some nontrivial constituent or relation inside it, from still more primitive explicit substrate structure in a way that changes benchmark-facing status. The benchmark boundary remains fixed throughout. The one operative metric is still exactly \(\CompletedMetric\), the ordinary matter route is unchanged, there is no second operative metric, the low-energy matter law is not enlarged, and no stronger promotion language is licensed. The positive result is therefore routing discipline below the body-response-core necessity / uniqueness frontier.
\end{abstract}

\section{Introduction}

The active derivational program of \TheoryName{} is no longer merely at pullback-body-response-core scope. The preceding necessity / uniqueness theorem already showed that, on the recorded section-centered hidden-response route, the current body-response core carries no remaining benchmark-facing presentational freedom beyond itself. That theorem therefore spent the smaller-collapse or smaller-uniqueness option that had remained open inside the prior body-response frontier.

The present note records the routing consequence of that gain. It does not reopen the larger pullback-body-operator, pullback-operator, pullback-quadratic, hidden-response, residual-normal-form, or pre-proto-section presentations as if they were again independently live. It records only which kind of next manuscript would now count as an honest benchmark-facing gain once the current burden sits on deriving or sharpening the already-unique body-response core itself.

\paragraph{Framework and claim boundary.}
\TheoryProgram{} denotes the broader \Aether{} / \Aflow{} framework built on the claims that the \Aether{} is the underlying four-dimensional substrate of reality, the \Aflow{} is its intrinsic ordered motion, observed three-dimensional space is the local experiential slice of that deeper substrate, S-time is the experienced order of change arising from matter, light, and the \Aflow{}, and observed expansion is the three-dimensional appearance of deeper four-dimensional ordered motion. Within that framework, \TheoryName{} denotes the exact-closure theory in which the effective gravitational dynamics are adopted to be exactly Einsteinian with universal matter coupling. In the active sequence, \emph{adoption} means use of that established relativistic dynamics without claiming substrate derivation, while \emph{derivation} is reserved for a first-principles recovery from explicit substrate variables.

The present note stays strictly inside that boundary. It records only which kind of next manuscript would now count as an honest benchmark-facing gain below the body-response-core necessity / uniqueness frontier.

\section{Fixed Inputs}

\begin{definition}[Body-response-core necessity / uniqueness frontier inputs]
The criterion recorded here uses the following fixed inputs.
\begin{enumerate}
    \item The benchmark exact-closure package, which fixes the public one-metric GR target of the repository.
    \item The benchmark gatekeeping layer, including the recorded safeguard that blocks same-output deeper-origin relays from counting as new front-line progress once the live benchmark-relevant datum is unchanged.
    \item The earlier theorem showing that the current pullback body-operator-core presentation \(\PullbackBodyOperatorCore\) carries no independent benchmark-facing freedom beyond the smaller observer-relevant pullback body-response core \(\PullbackBodyResponseCore\).
    \item The later theorem proving that the current observer-relevant pullback body-response core \(\PullbackBodyResponseCore\) is the unique minimal benchmark-facing datum on the recorded route up to body-response-faithful equivalence.
\end{enumerate}
\end{definition}

\begin{remark}
The present note does not replace those manuscripts. It records the routing consequence of reading them together under the active safeguard against recursive depth drift.
\end{remark}

\section{Why the Live Burden Now Sits Below Body-Response-Core Restatement}

\begin{proposition}[The live burden is renewed below body-response-core restatement]
At the current stage of the primitive-reservoir observer-reduction line, the body-response-core necessity / uniqueness theorem renews the live benchmark-facing burden below any mere restatement of the observer-relevant pullback body-response core \(\PullbackBodyResponseCore\).
\end{proposition}

\begin{proof}
That theorem changes benchmark-facing status because it proves that the present body-response core \(\PullbackBodyResponseCore\) is already the unique minimal benchmark-facing datum on the recorded route up to body-response-faithful equivalence. Once that uniqueness statement is in hand, another manuscript that merely preserves the same body-response core no longer removes any remaining benchmark-facing freedom. The live burden therefore no longer sits on the task of restating or repackaging \(\PullbackBodyResponseCore\) itself. It sits on deriving or sharpening that already-unique core from still more primitive explicit substrate structure.
\end{proof}

\begin{definition}[Benchmark-neutral same-core relay below the body-response necessity / uniqueness frontier]
A \emph{benchmark-neutral same-core relay below the body-response necessity / uniqueness frontier} is a manuscript that preserves the same benchmark one-metric output, the same ordinary matter route, the same current pullback body-response core \(\PullbackBodyResponseCore\) up to body-response-faithful equivalence, the same current pullback body-operator core \(\PullbackBodyOperatorCore\), the same current pullback-operator core \(\PullbackOperatorCore\), the same current pullback-quadratic core \(\PrePreProtoSectionPullbackCore\), the same current hidden-response package consequences, the same current residual-normal-form datum, the same current benchmark-faithful pre-proto-section package, and the same downstream benchmark-facing consequences while merely relocating the unexplained substrate layer deeper, renaming the same current core, or re-expanding the presentation back to the larger pullback-body-operator, pullback-operator, pullback-quadratic, hidden-response, residual-normal-form, or pre-proto-section presentations without a new derivation, new forcing theorem, or comparably sharp new gain.
\end{definition}

\section{The Current Post-Uniqueness Criterion}

\begin{definition}[Admissible next benchmark-facing manuscript below the body-response necessity / uniqueness frontier]
Below the current body-response-core necessity / uniqueness frontier, a manuscript counts as an admissible next benchmark-facing move only if it does at least one of the following:
\begin{enumerate}
    \item derives or justifies the current observer-relevant pullback body-response core \(\PullbackBodyResponseCore\) itself from still more primitive explicit substrate structure in a genuinely benchmark-facing way;
    \item derives or justifies some nontrivial constituent of \(\PullbackBodyResponseCore\) from still more primitive explicit substrate structure in a way that sharpens the present route;
    \item proves a sharper necessity, uniqueness, or compatibility theorem for some nontrivial constituent or relation inside \(\PullbackBodyResponseCore\);
    \item does one of the previous three while preserving the same benchmark one-metric package, the same ordinary matter route, and the same claim boundary.
\end{enumerate}
\end{definition}

\begin{theorem}[Next honest benchmark-facing task below the body-response necessity / uniqueness frontier]
The next honest benchmark-facing manuscript below the current body-response-core necessity / uniqueness frontier must derive or justify the current observer-relevant pullback body-response core \(\PullbackBodyResponseCore\), or some nontrivial constituent or relation inside it, from still more primitive explicit substrate structure in a benchmark-relevant way, or else prove a genuinely sharper theorem inside that same core. A manuscript that only repeats the same current body-response datum through another deeper relay, re-expands the discussion back to the larger pullback-body-operator, pullback-operator, pullback-quadratic, hidden-response, residual-normal-form, or pre-proto-section presentations without such a new gain, or invokes a quantum branch without doing the same benchmark-facing work does not satisfy the criterion.
\end{theorem}

\begin{proof}
By Proposition 1, the live burden already lies below any mere restatement of the current body-response core \(\PullbackBodyResponseCore\). Therefore a subsequent manuscript must improve on that already-unique core rather than merely accompany it. A deeper-origin manuscript can do so only if it changes benchmark-facing status by more than depth alone. If it simply reproduces the same current body-response core and its same downstream outputs, the routing safeguard classifies it as benchmark neutral. The remaining honest alternatives are therefore exactly those stated in the definition: derive or justify the body-response core itself, derive or justify some nontrivial constituent or relation inside it, or prove a sharper internal theorem there.
\end{proof}

\section{What the Next Move Should Not Be}

\begin{proposition}[Screened next moves below the body-response necessity / uniqueness frontier]
Below the current body-response-core necessity / uniqueness frontier, none of the following is the default next benchmark-facing move:
\begin{enumerate}
    \item another same-output deeper-origin relay that merely preserves the current pullback body-response core \(\PullbackBodyResponseCore\) up to body-response-faithful equivalence and therefore merely reproduces the same current pullback body-operator core \(\PullbackBodyOperatorCore\), the same current pullback-operator core \(\PullbackOperatorCore\), the same current pullback-quadratic core \(\PrePreProtoSectionPullbackCore\), the same current hidden-response package, the same current residual-normal-form datum, the same current benchmark-faithful pre-proto-section package, and the same downstream benchmark-facing consequences;
    \item another re-expansion back to the larger pullback-body-operator, pullback-operator, pullback-quadratic, hidden-response, residual-normal-form, or pre-proto-section presentations as if those already-spent larger presentations were again independently live burdens;
    \item a quantum branch that does not derive or justify the same body-response core, derive or justify a nontrivial constituent or relation inside it, or close a benchmark derivation gate in a genuinely new way;
    \item a manuscript that introduces a second operative metric, enlarges the ordinary matter law, or uses stronger promotion language.
\end{enumerate}
\end{proposition}

\begin{proof}
Item (1) fails the preceding criterion because it leaves the renewed live burden unchanged. Item (2) likewise fails because it re-expands to larger already-spent presentations rather than deriving or sharpening the already-unique body-response core. Item (3) does not directly address the live burden at current scope unless it performs the same benchmark-facing work named above. Item (4) violates the fixed claim boundary of the benchmark package and therefore cannot count as a disciplined next step on the active line. The benchmark package remains the exact one-metric GR package already fixed by adoption, so any admissible next manuscript must preserve that boundary.
\end{proof}

\begin{remark}
This screening statement is not a denial that those deeper or alternative manuscripts may matter strategically. It is a routing statement: they are not the default way to spend the next benchmark-facing move below the body-response necessity / uniqueness frontier unless they do the same benchmark-facing work named above.
\end{remark}

\section{Conclusion}

The positive result of this note is routing discipline below the body-response-core necessity / uniqueness frontier. The present necessity / uniqueness theorem remains the live benchmark-facing gain because it already proves that the current observer-relevant pullback body-response core is the unique minimal benchmark-facing datum on the recorded route. The earlier pullback-body-response-core collapse theorem remains preserved main-line history, but under the recursive-depth safeguard it is no longer itself the default current frontier.

The scope remains narrow and honest. The benchmark package is unchanged: the one operative metric is still exactly \(\CompletedMetric\), the ordinary matter route is unchanged, there is no second operative metric, and the low-energy matter law is not enlarged. Nothing here upgrades adoption to derivation or licenses stronger promotion language.

The next open burden is therefore clear. The next benchmark-facing manuscript should derive or justify the current body-response core itself, or some nontrivial constituent or relation inside it, from still more primitive explicit substrate structure in a way that changes benchmark-facing status. Another same-output deeper relay that merely preserves that body-response core, a re-expansion back to the larger pullback-body-operator, pullback-operator, pullback-quadratic, hidden-response, residual-normal-form, or pre-proto-section presentations, or a quantum branch that does not do that same benchmark-facing work is not the clean next manuscript target at current scope.

\end{document}
